
\begin{abstract}
We identify and represent the theta-semigroups derived from Liouville current
conditioning in Papers I--III.  The mechanical input determines the parameter:
$\theta=I_R'(j)$.  A separate rigidity theorem shows that, among smooth
law-invariant certainty equivalents, cash additivity and relevance leave only
the linear and exponential families; the nonzero mechanically selected field
chooses the entropic branch.

For the fixed-law physical diffusion, the terminal solution map is analytic on
bounded payoffs.  Its first derivative is expectation under a payoff-dependent
tangent probability and its second derivative is $\theta$ times covariance.
The tangent kernels satisfy a cocycle obtained by differentiating nonlinear
time consistency.  The microscopic Doob/Ruelle recursions carry exact tangent
path measures, and the deterministic homogenization of Paper III implies their
convergence to the diffusion tangent law.  A genuine stochastic-exponential
density gives Girsanov and quadratic/tangent BSDE representations.  Controlled
values are kept separate: their second variation contains the
optimizer-response correction from Paper IV.
\end{abstract}
\maketitle
\noindent\textbf{Platform identifier:} \texttt{TL1-LIO-v1}.

\tableofcontents

\section{The mechanically selected semigroup}\label{sec:mechanical}
Fix a radius $R$ and prepared current $j$ in the local LDP range.  Put
\[
 \theta=\theta_R(j)=I_R'(j),
 \qquad D=D_{R,\theta}.
\]
Let
\[
 dX_s=b(X_s)\,ds+\sqrt D\,dW_s
\]
be the physical comoving diffusion of Paper II.  For bounded terminal work
$\phi$ and running work $c$, define
\[
 \cE_{t,T}^{R,j}[\phi](x)
 =\frac1\theta\log\E_{t,x}
 \exp\left\{\theta\left[
 \phi(X_T)+\int_t^T c(X_s)\,ds\right]\right\},
\]
with the continuous linear limit at $\theta=0$.
Paper III proves that this is the limit of the canonical Liouville--Ruelle
excess-pressure recursion.

\section{A cash-additive rigidity theorem}\label{sec:rigidity}
Let $u:\R\to\R$ be strictly increasing and $C^2$, and define the certainty
equivalent
\[
 \mathfrak C_u(X)=u^{-1}(\E u(X))
\]
for bounded random variables on every finite probability space.

\begin{theorem}[Affine-or-exponential rigidity]\label{thm:rigidity}
Assume cash additivity
\[
 \mathfrak C_u(X+a)=\mathfrak C_u(X)+a
\]
for all bounded $X$ and constants $a$.  Then, after an affine change of $u$,
either
\[
 u(x)=x
\]
or
\[
 u(x)=e^{\gamma x}
\]
for one constant $\gamma\ne0$.  If the certainty equivalent is strictly
relevant and nonlinear, only the exponential branch remains.
\end{theorem}
\begin{proof}
Cash additivity implies that for every $a$, the map
$u(x)\mapsto u(x+a)$ preserves all convex combinations of values in the range
of $u$.  Hence there are scalars $A(a)>0$ and $B(a)$ with
\[
 u(x+a)=A(a)u(x)+B(a).
\]
Composition in $a$ gives
$A(a+b)=A(a)A(b)$.  Continuity yields either $A\equiv1$ or
$A(a)=e^{\gamma a}$.  In the first case, $u'$ is constant; in the second,
$u'$ is proportional to $e^{\gamma x}$.  Integrate and absorb affine
normalizations.  The linear branch is not strictly nonlinear.
\end{proof}

\begin{corollary}[Mechanical identification]\label{cor:identification}
Within the relevant nonlinear class, the only cash-additive certainty
equivalent compatible with the path-ensemble parameter selected in Paper I is
\[
 \frac1{\theta_R(j)}\log\E e^{\theta_R(j)X}.
\]
Thus mechanics supplies the numerical field and rigidity supplies the unique
functional form.
\end{corollary}

Dynamic law-invariant time-consistent monetary functionals satisfy a closely
related entropic rigidity under standard regularity hypotheses
\cite{KupperSchachermayer2009}.

\section{Analytic terminal maps and tangent probabilities}\label{sec:tangent}
For fixed $(t,x)$ and zero running cost, let
\[
 \mathfrak E(\phi)=\frac1\theta\log\E_{t,x}e^{\theta\phi(X_T)}.
\]

\begin{theorem}[Analytic tangent bundle]\label{thm:analytic}
The map $\mathfrak E:C_b(\R)\to\R$ is real analytic on bounded balls.  With
\[
 \frac{d\Qp^\phi}{d\Pp}
 =\frac{e^{\theta\phi(X_T)}}{\E e^{\theta\phi(X_T)}},
\]
one has
\[
 D\mathfrak E(\phi)\eta
 =\E_{\Qp^\phi}\eta(X_T),
\]
and
\[
 D^2\mathfrak E(\phi)(\eta,\zeta)
 =\theta\Cov_{\Qp^\phi}
